\documentclass[UTF8,a4paper,11pt]{ctexart}
\usepackage[margin=2.2cm]{geometry}
\usepackage{hyperref}
\usepackage{enumitem}
\setlist{nosep}
\hypersetup{colorlinks=true,linkcolor=black,urlcolor=blue}

\title{Proj57 明日操作手册}
\author{Proj57 初审材料}
\date{2026-06-30}

\begin{document}
\maketitle

\section{明天先准备什么}

建议准备：

\begin{itemize}
  \item Orange Pi 5 Plus 开发板。
  \item 稳定电源。
  \item HDMI 线和显示器。
  \item USB 键盘，最好再带鼠标。
  \item 笔记本电脑。
  \item 如果有，带一个 USB 转网口，但不是必须。
  \item 不要拔掉或重置板上的存储卡。
\end{itemize}

如果显示器能亮并进入系统，说明板子本身启动正常。若显示器无输出，再考虑使用 USB-TTL 串口模块看启动日志。

\section{第一步：确认板子是否启动}

接好电源、显示器、键盘后，看屏幕是否出现 Ubuntu、登录界面或终端。

如果能进入系统，打开终端执行：

\begin{verbatim}
uname -a
cat /etc/os-release
cat /proc/device-tree/model
free -h
df -h
\end{verbatim}

预期信息大致是：

\begin{verbatim}
Board: Orange Pi 5 Plus
OS: Ubuntu 22.04.5 LTS
Arch: aarch64
Memory: about 3.8 GiB
\end{verbatim}

\section{第二步：确认网络}

在板子终端执行：

\begin{verbatim}
ip addr
hostname -I
\end{verbatim}

如果看到类似下面的地址，说明板子在局域网里：

\begin{verbatim}
192.168.0.104
\end{verbatim}

然后在 Windows 笔记本 PowerShell 里执行：

\begin{verbatim}
ssh ubuntu@192.168.0.104
\end{verbatim}

如果 IP 变了，就把命令里的 IP 换成板子实际显示的 IP。

\section{第三步：确认 NPU 和 RKNN runtime}

在板子终端或 SSH 里执行：

\begin{verbatim}
ls -l /dev/dri/by-path | grep npu
ls -l /usr/lib/librknnrt.so
ls -l /usr/local/include/rknn
dmesg | grep -i rknpu | tail
\end{verbatim}

如果能看到 NPU 节点、\verb|librknnrt.so| 和 \verb|rknn_api.h|，说明板端 NPU runtime 基础条件存在。

\section{第四步：查看已经同步到板子的材料}

当前文档和小日志已经同步到：

\begin{verbatim}
/home/ubuntu/proj57_initial_review
\end{verbatim}

在板子上执行：

\begin{verbatim}
cd /home/ubuntu/proj57_initial_review
find . -maxdepth 3 -type f | sort
cat docs/initial_review/technical_doc.md
cat docs/initial_review/rknn_result.md
\end{verbatim}

这里面有初审技术文档、比赛背景、板端信息、NPU 探测结果、ONNX 正确性结果、性能表和演示脚本。

\section{第五步：本机继续跑模型验证}

在 Windows 的队伍仓库目录：

\begin{verbatim}
C:\Users\Laptop\Pictures\main\team-proj57
\end{verbatim}

使用本机 Python 环境可以执行：

\begin{verbatim}
..\.venv-ssh\Scripts\python.exe tools\check_onnx.py
\end{verbatim}

当前已经验证过的关键结果是：

\begin{verbatim}
frame_000000.jpg -> LEFT   expected: LEFT
frame_000227.jpg -> RIGHT  expected: RIGHT
\end{verbatim}

本机 ONNX Runtime CPU benchmark 约为：

\begin{verbatim}
frame_000000: 26.78 ms
frame_000227: 25.91 ms
\end{verbatim}

\section{第六步：RKNN 转换路线}

当前 Windows Python 不能直接安装 \verb|rknn-toolkit2|，所以不能直接在 Windows 环境完成 ONNX 到 RKNN 转换。正确路线是使用 x86\_64 Linux 环境，例如 WSL 或原生 Linux。

推荐命令如下：

\begin{verbatim}
cd /mnt/c/Users/Laptop/Pictures/main/team-proj57
python3 -m venv ~/.venvs/proj57-rknn
source ~/.venvs/proj57-rknn/bin/activate
python -m pip install -U pip
python -m pip install rknn-toolkit2==2.3.2
python -c "from rknn.api import RKNN; print('RKNN OK')"
python tools/convert_act_rknn.py
\end{verbatim}

当前已经成功生成：

\begin{verbatim}
build/act_rk3588_fp.rknn
\end{verbatim}

如果后续重转模型或更换量化策略失败，要保留完整日志，因为 ACT Transformer 结构仍可能遇到 RKNN 算子支持或内存问题。

\section{第七步：录制初审视频}

视频建议按这个顺序录：

\begin{enumerate}
  \item 展示比赛目标：ACT 模型适配 StarryOS、QEMU 和 RK3588。
  \item 展示 Orange Pi 5 Plus 实物和系统信息。
  \item 展示 NPU 节点、RKNN runtime 和内核日志。
  \item 展示任务三/QEMU/ONNX 的 LEFT/RIGHT 正确结果。
  \item 展示本地 ONNX benchmark。
  \item 展示 RKNN 转换和板端 runtime 已跑通的证据，以及下一步计划。
\end{enumerate}

重点不是把所有命令都讲复杂，而是让评委看懂：已经完成了模型链路验证、板端环境确认、RKNN/NPU 最小闭环和性能记录。

\section{遇到问题时先看这里}

\begin{itemize}
  \item 显示器没画面：先检查电源、HDMI 输入源、线缆，再考虑串口。
  \item SSH 不通：先在板子上用 \verb|hostname -I| 看 IP 是否变化。
  \item 没有网线口：优先用同一 Wi-Fi 或显示器键盘本机操作。
  \item RKNN 装不上：不要在 Windows Python 上硬装，切到 Linux/WSL。
  \item 不确定能否提交：先提交当前技术文档、ONNX 结果、RKNN 转换结果、板端 runtime 日志、性能表和演示视频；后续再冲刺 runtime 版本对齐和工程化。
\end{itemize}

\end{document}
